% Auto-generated; do not edit by hand.
\begin{proposition}[同步核的核版 Newman 阈值与 RH 相变点（resultant 消元，可审计）]
\label{prop:sync-kernel-weighted-newman-threshold}
\leanverified{paper\_sync\_kernel\_weighted\_newman\_threshold}
令 $B(u)=B_0+uB_1$ 为附录中的在线归一化同步核，并令其非零谱特征多项式为
\[
\begin{aligned}
P(\mu,u)&=\mu^6-(1+u)\mu^5-5u\mu^4+3u(1+u)\mu^3\\
&-u(u^2-3u+1)\mu^2+u(u^3-3u^2-3u+1)\mu+u^2(u^2+u+1).
\end{aligned}
\]
设 $\lambda_1(u)$ 为 Perron 根，$\Lambda(u)=\max_{i\ge2}|\lambda_i(u)|$。核版 RH 条件为 $\textsf{RH}_K(u):\Lambda(u)\le \lambda_1(u)^{1/2}$（定义 \ref{def:kernel-rh}）。
令 $b>0$，考虑联立方程 $P(-b,u)=0$ 与 $P(b^2,u)=0$，对 $b$ 消元得到结式分解
\[
\mathrm{Res}_b\!\bigl(P(-b,u),P(b^2,u)\bigr)=u^{9}(u-1)^{2}(u^{2}+u+1)(2u^{2}-u+3)\,Q(u),
\]
其中 $Q(u)\in\mathbb{Z}[u]$ 为次数 $37$ 的多项式：
\[
\begin{aligned}
Q(u)=u^{37} - 18 u^{36} + 123 u^{35}\\
&\ - 221 u^{34} - 1411 u^{33} + 4161 u^{32}\\
&\ + 22245 u^{31} - 52892 u^{30} - 158348 u^{29}\\
&\ - 42120 u^{28} + 606774 u^{27} + 3909659 u^{26}\\
&\ + 10387009 u^{25} - 37765093 u^{24} - 151950784 u^{23}\\
&\ + 25871406 u^{22} + 719435033 u^{21} + 601196630 u^{20}\\
&\ - 492437336 u^{19} + 626156178 u^{18} + 2164952510 u^{17}\\
&\ + 152029078 u^{16} - 730207162 u^{15} + 852147877 u^{14}\\
&\ + 184868548 u^{13} - 801264087 u^{12} - 85113751 u^{11}\\
&\ - 26931053 u^{10} - 93461933 u^{9} - 31696133 u^{8}\\
&\ - 1058443 u^{7} + 1222403 u^{6} + 137062 u^{5}\\
&\ - 10803 u^{4} - 2098 u^{3} - 61 u^{2}\\
&\ + 9 u + 1.
\end{aligned}
\]
在 $Q(u)=0$ 的 $(0,1)$ 内实根中，存在唯一根 $u_\star$ 使得
$|\mu_-(u_\star)|=\sqrt{\lambda_1(u_\star)}$（其中 $\mu_-(u)<0$ 为最坏非主谱半径所对应的负实根分支）。
该根给出同步核的 RH 相变点；并以 $u=e^{-\beta}$ 定义核版 Newman 阈值 $\beta_\star:=-\log u_\star$（定义 \ref{def:kernel-newman-threshold}）。
\[
\begin{aligned}
u_\star &\approx 0.7841489759097705,\\
\beta_\star=-\log u_\star &\approx 0.2431562563902237,\\
\lambda_1(u_\star) &\approx 2.66505879672116,\\
\mu_-(u_\star) &\approx -1.632500780006295,\\
\Lambda(u_\star) &\approx 1.632500780006295,\\
\Lambda(u_\star)/\sqrt{\lambda_1(u_\star)} &\approx 1.0,\\
\lambda_1(u_\star)-\mu_-(u_\star)^2 &\approx 0.0.
\end{aligned}
\]
数值核验表明：当 $0<u<u_\star$ 时 $\Lambda(u)<\lambda_1(u)^{1/2}$（严格 RH），当 $u=u_\star$ 时恰达界（RH-sharp），当 $u_\star<u\le1$ 时 $\Lambda(u)>\lambda_1(u)^{1/2}$（非 RH）。
\par\noindent\textbf{备注（审计护栏）：} $Q(u)$ 在 $(0,1)$ 内还存在额外实根（约为 0.136409571500976082）。但该根对应的是某个\emph{非 Perron}正实根与负实根之间的平方关系，并不触发 $\textsf{RH}_K$ 的达界。
\end{proposition}
